\chapter{Descripción del problema}

El dimensionamiento de un sistema híbrido eólico--fotovoltaico con electrólisis y almacenamiento electroquímico consiste en elegir las capacidades de sus componentes de modo que la energía útil producida a lo largo de un año resulte lo más barata posible, respetando las restricciones técnicas de operación y el límite regulatorio sobre la energía vendida a la red \cite{Khan2022Review,Li2024WPEB}. El criterio de optimización de esta tesis es único: el \textbf{costo nivelado de la energía (LCOE)}. El costo presente neto y el costo nivelado del hidrógeno según esta tesis se calculan para cada diseño, pero son indicadores secundarios que no intervienen en la comparación de soluciones.

Este capítulo enuncia el problema en su forma de optimización, define las ocho instancias sobre las que se resuelve, explica por qué es computacionalmente difícil y plantea la pregunta metodológica que los capítulos siguientes responden: con un presupuesto fijo y pequeño de evaluaciones, ¿qué estrategias de búsqueda, clásicas o híbridas, llegan al óptimo con mayor fiabilidad?

\section{Formulación del problema de dimensionamiento}
\label{sec:formulacion}

El sistema WPEB de Li et al. \cite{Li2024WPEB} instala una capacidad renovable total fija $C_{tot}$, repartida entre generación eólica $W$ y fotovoltaica $P=C_{tot}-W$; un electrolizador de capacidad $E$ que convierte electricidad en hidrógeno; y una batería de potencia $B^P$ y energía $B^E$. En la formulación de esta tesis el electrolizador se compra en $n_E$ módulos de $u$ MW, de modo que $E=u\,n_E$, y la batería se elige de un catálogo de potencias y duraciones, $B^E=B^P h_k$ con $h_k$ tomado de un conjunto finito $\{h_0,\ldots,h_{K-1}\}$. El vector de decisión es por tanto
\begin{equation}
\begin{aligned}
\mathbf{x}&=(W,\,n_E,\,B^P,\,k)\in\mathcal{X},\\
\mathcal{X}&=[0,C_{tot}]\times\{n_{min},\ldots,n_{max}\}\times\{0,\Delta_B,2\Delta_B,\ldots,B_{max}\}\times\{0,\ldots,K-1\},
\end{aligned}
\label{eq:dominio}
\end{equation}
con una coordenada continua y tres discretas. Cuando $B^P=0$ la duración carece de significado físico y se fija canónicamente en cero, de modo que distintos valores de $k$ describen el mismo diseño. El problema es
\begin{align}
\min_{\mathbf{x}\in\mathcal{X}}\quad & \mathrm{LCOE}(\mathbf{x})=\frac{C_{ann}(\mathbf{x})}{1000\,E_{use}(\mathbf{x})},
\qquad E_{use}=\sum_{t=1}^{8760}\left(P_{el,t}+P_{grid,t}\right)\Delta t \label{eq:problema}\\
\text{s.a.}\quad & P_{ren,t}+P_{desc,t}=P_{el,t}+P_{carga,t}+P_{grid,t}+P_{vert,t} \quad\forall t \tag{C1}\\
& W+P=C_{tot},\qquad W\geq0,\qquad P\geq0 \tag{C2}\\
& 0<E\leq0{,}5\,C_{tot} \tag{C3}\\
& 0\leq SOC_t\leq B^E,\qquad P_{carga,t},P_{desc,t}\leq B^P,\qquad P_{el,t}\in\{0\}\cup[\rho_{min}E,\,E] \quad\forall t \tag{C4}\\
& \mathrm{AGSR}(\mathbf{x})=\frac{E_{grid}}{E_{ren}}\leq\mathrm{AGSR}_{max} \tag{C5}
\end{align}
donde $C_{ann}$ es el costo anualizado del sistema (inversión, reemplazos y operación de todos los componentes, descontados y anualizados con el factor de recuperación de capital), $P_{el,t}$ y $P_{grid,t}$ son la potencia entregada al electrolizador y a la red en la hora $t$, $\rho_{min}$ es la carga mínima del electrolizador y AGSR es la fracción anual de energía renovable vendida a la red. El capítulo de marco teórico detalla los modelos de conversión, despacho y costos que vinculan $\mathbf{x}$ con esas magnitudes; el de diseño experimental precisa la procedencia y el tratamiento de cada restricción.

La propiedad que gobierna todo el trabajo es que $\mathrm{LCOE}(\mathbf{x})$ y $\mathrm{AGSR}(\mathbf{x})$ no tienen forma cerrada: son \textbf{salidas de un simulador horario} $\mathcal{S}(\mathbf{x};\text{datos})$ que recorre cronológicamente las 8.760 horas de un año meteorológico, despachando en cada hora la generación renovable entre electrolizador, batería, red y vertido según una prioridad fija. Las restricciones C1 y C4 las impone el propio despacho; C2 y C3 se garantizan por construcción del dominio; pero C5 sólo puede comprobarse al terminar el año. El problema pertenece así a la optimización de \textbf{caja negra} con una restricción verificable únicamente a posteriori.

\section{Instancias del problema}
\label{sec:instancias}

Una \textbf{instancia} del problema queda determinada por los datos meteorológicos y por la tupla de parámetros que define el dominio y la región factible,
\begin{equation}
\left(C_{tot},\;\mathrm{AGSR}_{max},\;u,\;n_{min},\;n_{max},\;\Delta_B,\;B_{max},\;\{h_k\}\right).
\end{equation}
Esta tesis resuelve ocho instancias sobre el mismo sitio y año, construidas a partir de la instancia de referencia cambiando el límite de ventas a red, la finura de los catálogos o la escala de la planta (Tabla~\ref{tab:instancias_auto}). Las cuatro primeras cambian un aspecto cada vez: \textbf{I1} es la instancia mixta de referencia; \textbf{I2} reduce a la mitad el límite de ventas a red, lo que estrecha la región factible; \textbf{I3} refina los catálogos del electrolizador y de la batería, lo que multiplica por cinco el número de combinaciones discretas sobre la misma región física; e \textbf{I4} escala la planta a 300 MW, lo que amplía los rangos de $W$ y $E$. Las cuatro restantes completan el cruce entre escala y restricción: \textbf{I5} reduce la planta a 100 MW; \textbf{I6} e \textbf{I7} combinan el límite de 0,10 con las escalas de 100 y 300 MW; e \textbf{I8} relaja el límite a 0,30, lo que ensancha la región factible. Las ocho comparten el simulador, la codificación de $\mathbf{x}$ y el tratamiento de restricciones, y su número responde a una necesidad del análisis: las preguntas sobre el desempeño \emph{en general} toman la instancia como unidad de análisis y requieren más de cuatro bloques para que un contraste tenga alguna potencia.

\TablaInstanciasAuto

Como $W$ es continuo, el dominio $\mathcal{X}$ no tiene cardinalidad finita. Sin embargo, al discretizar $W$ con paso de 0,01 MW cada instancia se convierte en un conjunto finito de entre \BFEDesigns{} y \BFCDesigns{} diseños anuales que esta tesis \textbf{enumera por completo}: el mejor diseño factible de esa enumeración es el \textbf{óptimo exacto del dominio discretizado} y constituye el único referente con el que se miden las brechas y tasas de éxito de los optimizadores. La distancia entre ese óptimo y el óptimo continuo queda acotada por el paso de discretización, y los optimizadores, que manipulan $W$ continuo, pueden quedar marginalmente por debajo de él.

\textbf{Enumeración y metaheurísticas.} La enumeración completa no vuelve innecesarios a los optimizadores: es el instrumento que permite medirlos. Es abordable en estas instancias por tres rasgos que no se conservan en extensiones realistas del problema: una única coordenada continua, catálogos discretos pequeños y un despacho por reglas fijas que puede evaluarse para muchos diseños a la vez. Aun así, exige una reimplementación vectorizada del despacho (Sección~\ref{sec:fuerza_bruta}), unas \VectorSpeedup{} veces más rápida por diseño que el simulador oficial, y entre \BruteForceMinutesMin{} y \BruteForceMinutesMax{} minutos de cómputo por instancia. Una sola coordenada continua adicional con la misma resolución, como la que resultaría de liberar la capacidad renovable total, multiplicaría por \SecondCoordFactor{} los diseños de I1 y llevaría su enumeración a unos \SecondCoordDays{} días con la misma implementación. Considerar varios años meteorológicos multiplicaría además el costo de cada evaluación, y dimensionar el almacenamiento de hidrógeno añadiría otra dimensión al dominio.

Una corrida metaheurística, en cambio, tiene un costo fijado por su presupuesto y no por el tamaño del dominio: \RequestsPerRun{} solicitudes, entre \DesignsPerBudgetMin{} y \DesignsPerBudgetMax{} veces menos que los diseños de cada instancia, que con el simulador oficial se resuelven en unos \RunSecondsMedian{} segundos. Las instancias se mantuvieron deliberadamente enumerables para que el óptimo fuera conocido y cada brecha midiera la distancia a un valor exacto, no a la mejor solución observada. Esta tesis no compara, por tanto, la enumeración con las metaheurísticas como métodos rivales: usa la primera para medir qué tan cerca del óptimo llegan las segundas con una fracción mínima del esfuerzo. Que ese desempeño se conserve en dominios donde la enumeración ya no es posible es una hipótesis que este trabajo no verifica.

\section{Naturaleza de la dificultad computacional}

\subsection{Evaluación costosa y sin forma analítica}

Determinar el LCOE de un diseño exige simular la operación del sistema durante un año completo con resolución horaria, porque tanto la producción de hidrógeno como la energía vendida a la red dependen del despacho hora a hora y no de magnitudes agregadas. La simulación es \textbf{secuencial}: el estado de carga de la batería en cada hora depende del estado precedente. En consecuencia, las solicitudes de función objetivo, y no el número nominal de iteraciones del algoritmo, constituyen la unidad del presupuesto de búsqueda. La ausencia de forma analítica impide además usar gradientes exactos, y la decodificación de las variables discretas crea mesetas en las que el gradiente aproximado es nulo. Otra vía sería formular el problema como un programa lineal entero mixto que optimice a la vez el tamaño y el despacho horario; esta tesis no la adopta porque el despacho es una regla de prioridad fija, la de Li et al., y optimizarlo cambiaría el problema que se reproduce y contra el que se miden los resultados.

\subsection{Región factible reducida, no convexa y verificable sólo a posteriori}

Las restricciones no acotan simplemente el dominio: definen una región factible de geometría irregular. En la malla base de Li et al., \BaseEvaluated{} puntos contienen sólo \BaseFeasible{} factibles (\BaseFeasiblePct\,\%); en las instancias mixtas enumeradas, la fracción factible del dominio discretizado depende sobre todo del límite de ventas a red:
va de \BFGFeasiblePct\,\% en I7 a \BFHFeasiblePct\,\% en I8, y las tres instancias con el límite de 0,10 no superan el \FeasibleShareTightMax\,\%, mientras que las cuatro con el límite de 0,20 se sitúan entre el \FeasibleShareMidMin{} y el \FeasibleShareMidMax\,\%. La restricción activa en el óptimo es C5: los mejores diseños se sitúan sobre el borde $\mathrm{AGSR}=\mathrm{AGSR}_{max}$, y como esa fracción es un resultado agregado del despacho, un algoritmo paga la simulación completa antes de saber si un candidato es admisible.

\subsection{Dominio mixto continuo--discreto}

Una coordenada continua y tres discretas obligan a decodificar cada propuesta antes de simularla: el número de módulos se redondea al entero más cercano, la potencia de batería al múltiplo del paso del catálogo y la duración al índice más próximo. Esa decodificación crea mesetas en el paisaje de búsqueda, pues propuestas distintas colapsan en el mismo diseño, y hace que los operadores de movimiento diseñados para espacios continuos gasten evaluaciones en diseños repetidos. La instancia I3, con catálogos más finos, reduce el tamaño de esas mesetas pero multiplica el número de combinaciones que un método debe discriminar.

\section{El problema metodológico y la pregunta de investigación}

A la dificultad computacional se superpone un problema de método. En la evaluación de algoritmos aleatorizados, las ejecuciones únicas, la ausencia de líneas base, la interpretación de medias sin contrastes ni tamaños de efecto y la confusión entre ejecuciones e instancias como unidad de análisis pueden conducir a conclusiones que el experimento no sostiene \cite{Demsar2006Statistical,Derrac2011Tutorial,Arcuri2014Hitchhiker}. Esta tesis no cuantifica la prevalencia de esas prácticas en la literatura de HRES-H$_2$; las trata como riesgos que el protocolo debe controlar.

La pregunta que organiza el trabajo es, entonces: \textbf{con un presupuesto idéntico y pequeño de \RequestsPerRun{} solicitudes por ejecución, ¿qué estrategias de búsqueda alcanzan el óptimo exacto de cada instancia con mayor fiabilidad, y aporta algo la hibridación PSO+GA respecto de los dos métodos clásicos de los que parte?} Para responderla se comparan siete algoritmos sobre las ocho instancias: cinco metaheurísticas clásicas, GWO, PSO, GA, SHO y WOA; un híbrido integrativo de bajo nivel, PSO+GA, que inyecta operadores genéticos en el enjambre de partículas; y la búsqueda aleatoria (RANDOM) como línea base sin aprendizaje. El híbrido conserva intacta la implementación del PSO y del GA de referencia y usa el mismo presupuesto. Su comparación con PSO aísla, en bloque, el efecto del componente genético añadido; su comparación con GA indica si esos mismos operadores rinden más dentro del enjambre que en el bucle generacional del algoritmo genético. La pregunta se descompone en tres, cada una con su unidad de análisis: si dos algoritmos difieren en una instancia dada, si difieren en general sobre el conjunto de instancias y si hay diferencias entre los siete a la vez. El protocolo incluye \Replicates{} ejecuciones independientes por algoritmo e instancia, hipótesis, nivel de significación y tolerancia práctica fijados a priori, contrastes no paramétricos corregidos por comparaciones múltiples, tamaños de efecto y análisis de sensibilidad frente a decisiones de modelado.
